\documentclass[a4paper]{article}
\usepackage{amsfonts}
\usepackage{amsmath}
\usepackage{amssymb}
\usepackage{graphicx}     

\begin{document}
  
\noindent \large Auteur: Dara van Engelen\\
\noindent \large Studierichting: Informatica\\
\noindent \large Oefeningenreeks 2D nr. 6.1\\

\medskip

\normalsize

$ f(x) = \dfrac{x^2 - x -6}{x^2 -2x-3}$  \\\\

$\lim\limits_{x\rightarrow -1} \left(\dfrac{x^2 - x -6}{x^2 -2x-3}\right) = \dfrac{-4}{0}$\\\\

$\lim\limits_{x\rightarrow -1-} \left(\dfrac{x^2 - x -6}{x^2 -2x-3}\right) = \lim\limits_{x\rightarrow -1-} \left(\dfrac{(x+2)(x-3)}{(x+1)(x-3)}\right)$ \\

$=\lim\limits_{x\rightarrow -1-} \left(\dfrac{x+2}{x+1}\right) =  -\infty$ \\\\

$\lim\limits_{x\rightarrow -1+} \left(\dfrac{x^2 - x -6}{x^2 -2x-3}\right) = \lim\limits_{x\rightarrow -1+} \left(\dfrac{(x+2)(x-3)}{(x+1)(x-3)}\right)$ \\

$=\lim\limits_{x\rightarrow -1+} \left(\dfrac{x+2}{x+1}\right) = +\infty$ \\\\

$\lim\limits_{x\rightarrow -2} \left(\dfrac{x^2 - x -6}{x^2 -2x-3}\right) = 0$ \\\\

$\lim\limits_{x\rightarrow 4} \left(\dfrac{x^2 - x -6}{x^2 -2x-3}\right) = \dfrac{6}{5}$ \\\\

$\lim\limits_{x\rightarrow 3} \left(\dfrac{x^2 - x -6}{x^2 -2x-3}\right) = \left(\dfrac{0}{0}\right) = \lim\limits_{x\rightarrow 3} \left(\dfrac{(x-3)(x-2)}{(x-3)(x+1)}\right) = \lim\limits_{x\rightarrow 3} \left(\dfrac{x+2}{x+1}\right) = \dfrac{5}{4}$ \\\\

$\lim\limits_{x\rightarrow \infty} \left(\dfrac{x^2 - x -6}{x^2 -2x-3}\right) = \lim\limits_{x\rightarrow \infty}\left(\dfrac{x^2}{x^2}\right) = 1$ \\

\begin{thebibliography}{999}
%\bibitem{a} Referentie
\end{thebibliography}
\end{document} 